\subsection{李代数的表现}

设 \(\mathfrak{g}\) 为李代数，\(\pmb{a} = (a_i)_{i\in I}\) 为 \(\mathfrak{g}\) 的元素族。设 \(f_{a}\) 为 \(\mathrm{L(I)}\) 到 \(\mathfrak{g}\) 的把每个 \(i\in I\) 映为 \(a_{i}\) 的同态。\(f_{a}\) 的像是 \(\mathfrak{g}\) 中由 \(\pmb{a}\) 生成的子代数；\(f_{a}\) 的核的元素称为族 \(\pmb{a}\) 的关系元。族 \(\pmb{a}\) 若 \(f_{a}\) 是满射（相应地，单射、双射），则称为生成的（相应地，自由的、基的）。

设 \(\mathfrak{g}\) 为李代数。\(\mathfrak{g}\) 的一个表现是指一个有序对 \((\boldsymbol{a},\boldsymbol{r})\)，由生成族 \(\boldsymbol{a} = (a_{i})_{i\in \mathbb{I}}\) 与关系元族 \(\boldsymbol{r} = (r_j)_{j\in \mathbb{J}}\)（\(\boldsymbol{a}\) 的关系元的族）组成，后者生成 L(I) 的理想即 \(f_{a}\) 的核。我们也说 \(\mathfrak{g}\) 由族 \(\boldsymbol{a}\) 连同关系元 \(r_j\)（\(j\in J\)）所表现。

设 I 为集合，\(\pmb{r} = (r_j)_{j \in J}\) 为自由李代数 L(I) 的元素族；设 \(a_r\) 为 L(I) 的由 \(\pmb{r}\) 生成的理想。商代数 L(I, \(\pmb{r}\) ) = L(I)/ \(a_r\) 称为由 I 与关系元族（\(r_j)_{j \in J}\)）所定义的李代数；我们也说 L(I, \(\pmb{r}\) ) 由表现 (I, \(\pmb{r}\) ) 定义，或由 (I; ( \(r_j = 0\) ) \(_{j \in J}\)) 定义。当族 \(\pmb{r}\) 为空时，L(I, \(\pmb{r}\) ) = L(I)。

设 \(I\) 与 \(r\) 如上；设 \(\xi_i\) 表示 \(i\) 在 \(L(I, r)\) 中的像。生成族 \(\xi = (\xi_i)_{i \in I}\) 与关系元族 \(r\) 构成 \(L(I, r)\) 的一个表现。反之，若 \(g\) 为李代数且 \((a, r)\)（其中 \(a = (a_i)_{i \in I}\)）是 \(g\) 的一个表现，则存在唯一的同构 \(u: L(I, r) \to g\)，使得 \(u(\xi_i) = a_i\) 对一切 \(i \in I\) 成立。

